\documentclass[10pt,a4paper]{amsart}
\usepackage[margin=19mm]{geometry}
\usepackage{amsmath,amssymb,mathtools}
\usepackage[scr=rsfs,cal=euler]{mathalfa}
\usepackage{xfrac}
\usepackage[unicode,hidelinks,bookmarks=false]{hyperref}
\newcommand{\ie}{i.e.\ }
\newcommand{\cf}{cf.\ }
\newcommand{\resp}{respectively\ }
\newcommand{\Subsection}{\subsection}
\providecommand{\nobreakdash}{\nobreak\hbox{-}}
\setlength{\emergencystretch}{2em}
\multlinegap0pt
\usepackage{bm}
\usepackage{bbm}
\usepackage{dsfont}
\usepackage{xspace}
\usepackage{cancel}
\usepackage{mathtools}
\usepackage{amsthm}
\makeatletter
\@ifundefined{theorem}{\newtheorem{theorem}{Theorem}}{}
\@ifundefined{prop}{\newtheorem{prop}[theorem]{Proposition}}{}
\@ifundefined{remark}{\newtheorem{remark}[theorem]{Remark}}{}
\makeatother
\title[Moment Lagrangians, unobstructedness and symplectic groupoid]{Moment Lagrangians, unobstructedness and symplectic groupoids}
\author{Yan-Lung Leon Li}
\date{Source: arXiv:2609.12033v1 (CC BY 4.0)}
\begin{document}
\maketitle

\noindent\emph{Source and licence.} Moment Lagrangians, unobstructedness and symplectic groupoids. Version arXiv:2609.12033v1 (CC BY 4.0), distributed under the Creative Commons Attribution 4.0 International licence, as recorded in the arXiv licence metadata for this exact version and shown on the abstract page https://arxiv.org/abs/2609.12033v1. Full rights notice: licenses/arxiv-2609.12033-CC-BY-4.0.txt.\par\smallskip

\noindent\textit{Editorial scope.} Counted unit: the Remark whose source label is \texttt{None} in the article (source0.tex lines 388--390, source label \texttt{None}), delivered together with the article's own complete original proof of it (source0.tex lines 392--420, one \texttt{proof} environment of 29 lines). No mathematics is written, repaired or re-derived: statement and proof are the article's own text line for line. Required context retained uncounted: hamilton (lines 295--297); lagrangiansubmanifold (lines 355--362); twisted-groupoid (lines 370--384); symplecticform (lines 541--543). Every definition, display and label of those lines is retained unchanged. Omitted from this package and not redistributed: 1--294, 298--354, 363--369, 385--387, 391--391, 421--540, 544--577. Nothing inside a retained excerpt is omitted or altered: every retained body is delivered exactly as the article's lines are written, so no omission receipt of any kind accompanies this package. Citations: the retained excerpts carry no citation command at all, so no bibliography entry of the article is transcribed and no reference is redistributed. Reference adaptation: none was needed and none was made; every reference inside the delivered unit resolves inside it, so no unresolved reference or citation is produced. Printed slips kept literal: no printed word, symbol, hypothesis or number of the counted statement or of its proof was altered. Layout and class: the article's own class is \texttt{amsart} with packages including xcolor, amsfonts, amssymb, graphicx, mathrsfs, amscd, hyperref, stmaryrd, caption, subcaption, color, comment, bm, bbm; the wrapper is the lane's pinned \texttt{amsart} editorial wrapper at 10pt on A4, which restates the theorem-like environments the retained text needs and adds the article's own macro definitions that the retained text needs (none), with extra wrapper packages (bm, bbm, dsfont, xspace, cancel, mathtools). The delivered numbering is the unit's own. Assets: no retained span contains a figure environment, a graphics-inclusion command, a PDF-page inclusion, a table or a TikZ picture; no image file is copied into this package, the package declares no asset dependency and no image file is redistributed with it. Licence: version arXiv:2609.12033v1 of this article is distributed under the Creative Commons Attribution 4.0 International licence, as recorded in the arXiv licence metadata for this exact version and shown on the abstract page https://arxiv.org/abs/2609.12033v1. A full rights notice is delivered at licenses/arxiv-2609.12033-CC-BY-4.0.txt. Characters: every retained excerpt is pure ASCII, so the delivered engine, XeLaTeX with its default Latin Modern text font, reports no missing character.
\begin{equation}\label{hamilton}
  \iota_{v^\#}\omega = \langle d\mu(-), v \rangle, \forall v\in \mathfrak{g}.
\end{equation}

\begin{prop} \label{lagrangiansubmanifold}
Given a symplectic groupoid $(X, \omega) \xRightarrow[t]{s} M$, the following are Lagrangians:
\begin{enumerate}
\item $\underline{m} := Graph(m) = \{(p,q,p\cdot q) \ |\,  (p,q)\in X^{(2)}\} \subseteq X^- \times X^- \times X,$
\item $\underline{e} := e(M)=\{e(x) |\,  x\in M\} \subseteq X,$
\item $\underline{i} := Graph(i) =\{(p,p^{-1})) \in X \times X|\,  p\in X \}\subseteq X \times X.$ 
\end{enumerate}
\end{prop} 

\begin{theorem}\label{twisted-groupoid}
    Given a Hamiltonian $G$-space $(Y, \omega, \mu)$, there exists a symplectic groupoid structure on $(X \coloneqq (G \times \mathfrak{g}^*)^- \times Y^- \times Y, \omega_X \coloneqq ((-\Omega)\oplus (-\omega) \oplus \omega))$:
    $$((X , \omega_X) \xRightarrow[t]{s} G \times Y, m, e, i),$$
given by the following:
    \begin{itemize}
        \item  $s(g, \xi, y, z) = (g, g^{-1}\cdot z)$; $t(g, \xi, y, z) = (g, y)$;
        \item $(g_1, \xi_1, y_1, z_1) \cdot (g_2, \xi_2, y_2, z_2) = (g, \xi_1 + \xi_2 - \mu(y_2), y_1, z_2)$, where 
        $$g_1 = g_2 = g; z_1 = g\cdot y_2;$$
        \item $e(g, x) = (g, \mu(x), x, g\cdot x)$;
        \item $i(g, \xi, y, z) = (g, \mu(g^{-1}\cdot z)+ \mu(y) - \xi, g^{-1}\cdot z, g\cdot y)$.
    \end{itemize}

     In particular, $L_\mu$ is the unit and also the fixed locus of the anti-symplectic involution $i$:
    $$L_\mu = \underline{e} = Fix(i).$$
\end{theorem}

\begin{align}\label{symplecticform}
\Omega((Z, \eta), (W, \zeta)) = \langle \eta, \mathcal{L}^{-1}_{g}(W) \rangle - \langle \zeta, \mathcal{L}^{-1}_{g}(Z) \rangle - \langle \xi, [\mathcal{L}^{-1}_{g}(Z), \mathcal{L}^{-1}_{g}(W)] \rangle.
\end{align}

\begin{remark}
Combining Theorem \ref{twisted-groupoid} with Proposition \ref{lagrangiansubmanifold}, we obtain an alternative proof of the fact that the moment Lagrangian $L_\mu$ is a Lagrangian of $(X, \omega_X)$.
\end{remark}

\begin{proof}
It is straightforward to check that the above maps satisfy the axioms of a Lie groupoid, and that $L_\mu = \underline{e} = Fix(i)$. It remains to check that $\omega_X$ is multiplicative, which will be done by definition as follows.

    First of all, the tangent space of $X^{(2)}$ at $(x_1, x_2) \coloneqq ((g, \xi_1, y_1, z_1), (g, \xi_2, y_2, z_2))$ is given by 
    $$T_{(x_1, x_2)}X^{(2)} = \{((Z, \eta_1, v_1, v'_1), (Z, \eta_2, v_2, v'_2))| Z\in \mathfrak{g}; v'_1 = g_* v_2 + (y_2)_* Z\}.$$
    where $g_*: T_{y}Y \rightarrow T_{g\cdot y}Y$ is the differential of the action by $g$ and $y_*: \mathfrak{g} \rightarrow T_{y}Y$ is the infinitesimal generator $Z^\#$ evaluated at $y$.

    Also, the differential of $m$ at $(x_1, x_2)$ is given by
 $$dm_{(x_1, x_2)}((Z, \eta_1, v_1, v'_1), (Z, \eta_2, v_2, v'_2))=(Z, \eta_1 + \eta_2 - d\mu(v_2), v_1, v'_2).$$

    Therefore, the symplectic form is multiplicative if and only if for any 
    $$((Z, \eta_1, v_1, v'_1), (Z, \eta_2, v_2, v'_2)), (W, \xi_1, w_1, w'_1), (W, \xi_2, w_2, w'_2))\in T_{(x_1, x_2)}X^{(2)},$$
    the following holds:
\begin{align*}
    &\omega_X((Z, \eta_1 + \eta_2 - d\mu(v_2), v_1, v'_2), (W, \xi_1 + \xi_2 - d\mu(w_2), w_1, w'_2))\\
    = &\omega_X((Z, \eta_1, v_1, v'_1), (W, \xi_1, w_1, w'_1)) + \omega_X((Z, \eta_2, v_2, v'_2), (W, \xi_2, w_2, w'_2)).
\end{align*}
Using \ref{symplecticform}, the above is equivalent to the following:
\begin{align*}
&-\langle d\mu(v_2), \mathcal{L}^{-1}_{g}(W) \rangle + \langle d\mu(w_2), \mathcal{L}^{-1}_{g}(Z) \rangle + \langle \mu(y_2), [\mathcal{L}^{-1}_{g}(Z), \mathcal{L}^{-1}_{g}(W)] \rangle \\
= &-\omega(g_*(v_2), (y_2)_*W) - \omega((y_2)_*Z, g_*(w_2)) - \omega((y_2)_*Z, (y_2)_*W).
\end{align*}
The above equality holds if the following identities hold:
\begin{align*}
\langle d\mu(v), \mathcal{L}^{-1}_{g}(W) \rangle &= \omega(g_*(v), (y_2)_*W);\\
\langle \mu(y_2), [\mathcal{L}^{-1}_{g}(Z), \mathcal{L}^{-1}_{g}(W)] \rangle &= - \omega((y_2)_*Z, (y_2)_*W).
\end{align*}
The first identity follows from the Hamilton's equation \ref{hamilton}, and the second identity follows from that the comoment map $\mu^*$ is a Lie algebra homomorphism.
\end{proof}

\end{document}
